% 附录：\appendix 由 main.tex 发出，附录页码独立计数，
% 页码切换为 A-当前 / A-总数（见 inkline/layout.tex）。
% 不需要附录时，把 main.tex 里的 \appendix 与该文件的 \input 一起删除即可。
\section{附录}

\begin{frame}[allowframebreaks]{参考文献}
  % 文献表字体由 inkline/base.tex 的 \bibfont 统一设为 Times 风格衬线体，
  % 中文条目仍由 ctex 回退到思源宋体，无需在此另行设置。
  \small
  \printbibliography[heading=none]
\end{frame}

% 术语与符号表：首次出现的缩写或记号可在此统一解释，避免正文反复展开
\begin{frame}{术语与符号}
  \small
  \begin{tabularx}{\linewidth}{@{}lY@{}}
    \toprule
    \textbf{符号 / 缩写} & \textbf{含义} \\
    \midrule
    符号 A & 在此给出定义与单位 \\
    缩写 B & 全称及首次出现位置 \\
    缩写 C & 全称及首次出现位置 \\
    \bottomrule
  \end{tabularx}
  \inksource{符号表按正文出现顺序排列；正式使用时替换为实际记号。}
\end{frame}

% framebreak 手动分页；内容过长时 allowframebreaks 也可自动分页。
\begin{frame}[allowframebreaks]{长内容续页}
  \textbf{第一页}
  \begin{itemize}
    \item allowframebreaks 允许一个 frame 分成多页。
    \item framebreak 可在需要的位置手动分页。
  \end{itemize}
  \inknote{续页标题会追加序号。此示例放在附录，以便同时查看附录页码。}
  \framebreak
  \textbf{第二页}
  \begin{itemize}
    \item 长文献列表可按相同方式自动分页。
    \item 分步动画请使用普通 frame，避免与自动分页混用。
  \end{itemize}
\end{frame}

\begin{frame}[plain,c,noframenumbering]
  \centering
  {\usebeamerfont{title}\usebeamercolor[fg]{title}示例结束\par}
  \vspace{0.8em}
  普通汇报可在这里放置结束语。
  \inknote{plain 隐藏页眉页脚；c 让内容垂直居中；
  noframenumbering 使此页不占用页码。}
\end{frame}
